\enquestion{7.21}{
    Let $X_{1}, X_{2}, \ldots, X_{7}$ denote a random sample from a population having mean $\mu$ and variance $\sigma^{2}$.
    Consider the following estimators of $\mu$:
    \[
        \hat{\Theta}_{1} = \frac{X_{1} + X_{2} + \dots + X_{7}}{7}
    \]
    \[
        \hat{\Theta}_{2} = \frac{2X_{1} - X_{6} + X_{4}}{2}
    \]
    \begin{enumerate}[label=(\alph*)]
        \item Is either estimator unbiased?
        \item Which estimator is better?
        In what sense is it better?
        Calculate the relative efficiency of the two estimators.
    \end{enumerate}
}
\enanswer{
    \begin{enumerate}[label=(\alph*)]
        \item Both estimators are unbiased because $E(\hat{\Theta}_{1}) = \mu$ and $E(\hat{\Theta}_{2}) = \mu$.
        \item $\hat{\Theta}_{1}$ is better because it is more precise. Since $V(\hat{\Theta}_{1}) = \frac{\sigma^{2}}{7}$ and $V(\hat{\Theta}_{2}) = \frac{3\sigma^{2}}{2}$, the relative efficiency of $\hat{\Theta}_{2}$ to $\hat{\Theta}_{1}$ is $\frac{2}{21}$.
    \end{enumerate}
}

\enquestion{7.23}{
    Suppose that $\hat{\Theta}_{1}$ and $\hat{\Theta}_{2}$ are estimators of the parameter $\theta$.
    We know that $E(\hat{\Theta}_{1}) = \theta$, $E(\hat{\Theta}_{2}) = \theta / 2$, $V(\hat{\Theta}_{1}) = 10$, $V(\hat{\Theta}_{2}) = 5$.
    Which estimator is better?
    In what sense is it better?
}
\enanswer{
    $\hat{\Theta}_{1}$ is preferred. $\hat{\Theta}_{1}$ is unbiased with $V(\hat{\Theta}_{1}) = 10$. $\hat{\Theta}_{2}$ is biased, but $2\hat{\Theta}_{2}$ is unbiased with $V(2\hat{\Theta}_{2}) = 20$. Since $\hat{\Theta}_{1}$ has a smaller variance among the unbiased options, it is the better estimator.
}

\enquestion{7.25}{
    Let three random samples of sizes $n_{1} = 8$, $n_{2} = 14$, and $n_{3} = 6$ be taken from a population with mean $\mu$ and variance $\sigma^{2}$.
    Let $S_{1}^{2}$, $S_{2}^{2}$, and $S_{3}^{2}$ be the sample variances.
    Show that $S^{2} = (7S_{1}^{2} + 13S_{2}^{2} + 5S_{3}^{2}) / 25$ is an unbiased estimator of $\sigma^{2}$.
}
\enanswer{
    Since $E(S_{1}^{2}) = \sigma^{2}$, $E(S_{2}^{2}) = \sigma^{2}$, and $E(S_{3}^{2}) = \sigma^{2}$, we have:
    \[
        E(S^{2}) = \frac{7}{25}\sigma^{2} + \frac{13}{25}\sigma^{2} + \frac{5}{25}\sigma^{2} = \sigma^{2}
    \]
    Thus, $S^{2}$ is an unbiased estimator of $\sigma^{2}$.
}

\enquestion{7.29}{
    Data on the oxide thickness of semiconductor wafers are as follows: $425$, $431$, $416$, $420$, $421$, $436$, $418$, $410$, $431$, $433$, $423$, $426$, $411$, $435$, $436$, $428$, $411$, $426$, $407$, $437$, $422$, $428$, $413$, $416$.
    \begin{enumerate}[label=(\alph*)]
        \item Calculate a point estimate of the mean oxide thickness for all wafers in the population.
        \item Calculate a point estimate of the standard deviation of oxide thickness for all wafers in the population.
        \item Calculate the standard error of the point estimate from part (a).
        \item Calculate a point estimate of the median oxide thickness for all wafers in the population.
        \item Calculate a point estimate of the proportion of wafers in the population that have oxide thickness of more than $430$ angstroms.
    \end{enumerate}
}
\enanswer{
    \begin{tabular}{*{3}{m{.3\textwidth}}}
        (a) $423.3333$ & (b) $9.1493$ & (c) $1.8676$ \\
        (d) $424$ & (e) $0.2917$ & \\
    \end{tabular}
}

\enquestion{7.32}{
    Two different plasma etchers in a semiconductor factory have the same mean etch rate $\mu$.
    However, machine $1$ is newer than machine $2$ and consequently has smaller variability in etch rate.
    We know that the variance of etch rate for machine $1$ is $\sigma_{1}^{2}$, and for machine $2$, it is $\sigma_{2}^{2} = a\sigma_{1}^{2}$.
    Suppose that we have $n_{1}$ independent observations on etch rate from machine $1$ and $n_{2}$ independent observations on etch rate from machine $2$.
    \begin{enumerate}[label=(\alph*)]
        \item Show that $\hat{\mu} = \alpha \bar{X}_{1} + (1 - \alpha) \bar{X}_{2}$ is an unbiased estimator of $\mu$ for any value of $\alpha$ between zero and one.
        \item Find the standard error of the point estimate of $\mu$ in part (a).
        \item What value of $\alpha$ would minimize the standard error of the point estimate of $\mu$?
        \item Suppose that $a = 4$ and $n_{1} = 0.5 n_{2}$.
        What value of $\alpha$ would you select to minimize the standard error of the point estimate of $\mu$?
        How \enquote{bad} would it be to arbitrarily choose $\alpha = 0.5$ in this case?
    \end{enumerate}
}
\enanswer{
    \begin{enumerate}[label=(\alph*)]
        \item $E(\hat{\mu}) = \alpha E(\bar{X}_{1}) + (1 - \alpha) E(\bar{X}_{2}) = \alpha\mu + (1-\alpha)\mu = \mu$
        \item $\text{S.E.} = \sigma_{1}\sqrt{\frac{\alpha^{2}}{n_{1}} + \frac{a(1-\alpha)^{2}}{n_{2}}}$
        \item $\alpha = \frac{an_{1}}{n_{2} + an_{1}}$
        \item $\alpha = \frac{2}{3}$, choosing $\alpha = 0.5$ increases the standard error and is thus not preferable.
    \end{enumerate}
}